\emph{Setting.} As in Theorem \claimref{synthesis.tower-top-pair-forgets-reversal}, $L_{p,q}$ is the union of the floors, spanned by the finite blades in the ground and plane generators, and conjugation by a generator $g$ acts diagonally on blades: for a blade $B$, $gB = (-1)^{|B|}Bg$ if $g \notin B$ and $gB = (-1)^{|B|-1}Bg$ if $g \in B$ (the blade argument of Theorem \claimref{synthesis.plane-commutant-even-ground}, Step 1). Since $\alpha_{B}$ is an automorphism of $B_{\infty}$ and the plane generators generate $B_{\infty}$, $x \in Z^{1}$ if and only if $xg = -gx$ for every plane generator $g$: if $x$ anticommutes with $g$ and $h$ then $x\,gh = gh\,x = \alpha_{B}(gh)x$, and so on for longer products.

\emph{(i).} Because conjugation by each generator is diagonal on blades, $x \in Z^{1}$ if and only if every blade in its support anticommutes with every plane generator. Let $B$ be such a blade. A blade involves finitely many planes, so some plane generator $h \notin B$, and $hB = -Bh$ forces $|B|$ odd. If some plane generator $g$ lay in $B$, then $gB = (-1)^{|B|-1}Bg = Bg$, a contradiction. So $B$ is an odd blade in the ground generators alone; conversely every such blade anticommutes with every plane generator. Hence $Z^{1} = \Cl^{1}(p,q)$. With $Z^{0} = \Cl^{0}(p,q)$ (Theorem \claimref{synthesis.plane-commutant-even-ground}(i)), $Z^{0} \oplus Z^{1}$ is the span of all ground blades, the subalgebra $\Cl(p,q) = T_{0}$, and the decomposition is its grading. An isomorphism of triples $\Psi$ (with $\Psi(B_{\infty}) = B_{\infty}$, $\Psi\alpha_{B} = \alpha'_{B}\Psi$) maps $Z^{0}$ onto $Z'^{0}$, since $\Psi(B_{\infty}) = B_{\infty}$, and $Z^{1}$ onto $Z'^{1}$: if $xb = \alpha_{B}(b)x$ for all $b$, then $\Psi(x)\Psi(b) = \Psi(\alpha_{B}(b))\Psi(x) = \alpha'_{B}(\Psi(b))\Psi(x)$, and $\Psi(b)$ ranges over $B_{\infty}$; the same with $\Psi^{-1}$ gives equality. So $\Psi$ restricts to a grading-preserving isomorphism $\Cl(p,q) \to \Cl(p',q')$.

\emph{(ii).} If such a $\Psi$ exists, its restriction to the commutants gives $\Cl^{0}(p,q) \cong \Cl^{0}(p',q')$, so $n' = n$ and $r' \equiv \pm r \pmod 8$ by Theorem \claimref{synthesis.plane-commutant-even-ground}(ii), and by (i) $\Cl(p,q) \cong \Cl(p',q')$. If $r' \equiv -r$ with $r \not\equiv 0, 4$, then $\Cl(p',q') \cong \Cl(q,p)$ (the real Clifford algebra depends only on $n$ and the residue modulo $8$), and $\Cl(p,q) \not\cong \Cl(q,p)$: the types for $r$ and $-r$ are $\R \oplus \R$ against $\C$ ($r = \pm 1$), $\R$ against $\HH$ ($r = \pm 2$) and $\C$ against $\HH \oplus \HH$ ($r = \pm 3$), Theorem \claimref{transport.classification-real-clifford-algebras}. So $r' \equiv r$. An algebra invariant that exhibits the difference without the classification is the signature of the regular trace form $(x,y) \mapsto \operatorname{tr}(L_{xy}) = 2^{n}\br{xy}_{0}$, diagonal on blades with entries $2^{n}B^{2}$; the check computes it for all $n \le 8$ and finds it different for $\Cl(p,q)$ and $\Cl(q,p)$ exactly when $r \not\equiv 0, 4 \pmod 8$. Conversely, for $n' = n$ and $r' \equiv r$, $p' \equiv p \pmod 4$, and the four-generator shift $x_{i} \mapsto x_{i}\omega_{4}$ of Theorem \claimref{synthesis.tower-top-pair-forgets-reversal}(ii), repeated, is an isomorphism fixing every plane generator, hence $B_{\infty}$ pointwise, hence respecting $\alpha_{B}$. The reversal $\Phi$ of Theorem \claimref{synthesis.tower-top-pair-forgets-reversal}(i) has $\Phi(e_{2}) = e_{1}f_{1}$ and $\Phi(f_{1}) = e_{1}e_{2}$, even elements, while $\alpha_{B}$ negates $e_{2}$ and $f_{1}$ and fixes $e_{1}f_{1}$, $e_{1}e_{2}$; so $\Phi\alpha_{B} \ne \alpha'_{B}\Phi$.

\emph{(iii).} An element $a + b\,e_{j} + c\,f_{j} + d\,e_{j}f_{j}$ of $P_{j}$ commutes with $e_{k}$ and $f_{k}$, $k \ne j$, if and only if $b = c = 0$, since $e_{j}, f_{j}$ anticommute with them and $1, e_{j}f_{j}$ commute. So, writing $C(\cdot)$ for the commutant in the limit, $P_{j} \cap C(P_{k}) = \mathrm{span}\{1, e_{j}f_{j}\} \cong \R \oplus \R$ ($(e_{j}f_{j})^{2} = -e_{j}^{2}f_{j}^{2} = 1$), whose elements of square $1$ are $\pm 1$ and $\pm e_{j}f_{j}$ (from $(a + d\,e_{j}f_{j})^{2} = a^{2} + d^{2} + 2ad\,e_{j}f_{j}$). Conjugation by $u = e_{j}f_{j}$ ($u^{-1} = u$) negates $e_{j}$ and $f_{j}$ and fixes $1$ and $u$, so $\alpha_{B}|_{P_{j}} = \operatorname{Ad}(u)$. Let $\Psi$ be a pair isomorphism with $\Psi(P_{j}) = P'_{\pi(j)}$, $\pi$ injective. For $j$ pick $k \ne j$; then $\Psi(P_{j} \cap C(P_{k})) = P'_{\pi(j)} \cap C(P'_{\pi(k)})$, the span of $1$ and $u' = e'_{\pi(j)}f'_{\pi(j)}$, so $\Psi(u) = \pm u'$ and, for $b \in P_{j}$, $\Psi(\alpha_{B}(b)) = \Psi(u)\Psi(b)\Psi(u)^{-1} = \operatorname{Ad}(u')(\Psi(b)) = \alpha'_{B}(\Psi(b))$. The $P_{j}$ generate $B_{\infty}$, so $\Psi\alpha_{B} = \alpha'_{B}\Psi$ on $B_{\infty}$, and (ii) applies. The shift fixes every plane, so the framed classification is that of (ii).

\emph{(iv), general clocks.} Write $\alpha$ for $\alpha'_{B}$ on the target and $L' = L_{p',q'}$, with ground generators $y_{1}, \dots, y_{n'}$. Three facts about $B_{\infty}$ are used. (a) It is simple with centre $\R$: it is the union of the floors $B_{k} \cong \Cl(k,k) \cong \Mat_{2^{k}}(\R)$ under unital embeddings, an element central in $B_{\infty}$ is central in some $B_{k+1}$ and so a scalar, and a nonzero two-sided ideal meets some $B_{k}$ in a nonzero ideal, hence contains $1$. (b) $\alpha$ is not inner on $B_{\infty}$: an implementer $d \in B_{k}$ would commute with $e_{k+1}\Omega_{k}$, $\Omega_{k} = \prod_{j \le k}e_{j}f_{j}$, which commutes with $B_{k}$ and is odd. (c) $L' = \bigoplus_{B}y_{B}B_{\infty}$ over the ground blades $y_{B}$, uniquely, and $b\,y_{B} = y_{B}\,\alpha^{|B|}(b)$ for $b \in B_{\infty}$.

Let $\beta = \Psi|_{B_{\infty}}$, an automorphism of $B_{\infty}$, $\tilde\alpha = \beta\alpha_{B}\beta^{-1}$, and $z_{i} = \Psi(x_{i})$. From $x_{i}g = -gx_{i}$ for the plane generators, $z_{i}b = \tilde\alpha(b)z_{i}$ for all $b \in B_{\infty}$. Put $w = z_{1}y_{1}$, invertible; since $y_{1}b = \alpha(b)y_{1}$, $wb = \gamma(b)w$ with $\gamma = \tilde\alpha\alpha$. Expand $w = \sum_{B}y_{B}b_{B}$ by (c): comparing components, $b_{B}b = \alpha^{|B|}(\gamma(b))\,b_{B}$ for every $B$ and $b$. For $b_{B} \ne 0$ the set $B_{\infty}b_{B} = b_{B}B_{\infty}$ is a nonzero two-sided ideal, so $b_{B}$ is invertible by (a) and $\alpha^{|B|}\gamma = \operatorname{Ad}(b_{B})$. If $|B|$ were odd, $\tilde\alpha = \gamma\alpha = \alpha\operatorname{Ad}(b_{B})\alpha = \operatorname{Ad}(\alpha(b_{B}))$ would be inner, and so would $\alpha_{B} = \beta^{-1}\tilde\alpha\beta$, against (b) for the source. So $|B|$ is even, $\gamma = \operatorname{Ad}(c)$ with $c = b_{B}$, and $\tilde\alpha = \operatorname{Ad}(c)\alpha$. Since $\tilde\alpha^{2} = 1$, $\operatorname{Ad}(c\,\alpha(c)) = 1$ and $c\,\alpha(c) = \mu \in \R \setminus \{0\}$ by (a). Then $R_{i} := c^{-1}z_{i}$ satisfies $R_{i}b = \alpha(b)R_{i}$ for all $b \in B_{\infty}$, so $R_{i} \in \Cl^{1}(p',q')$ by (i), and $z_{i} = cR_{i}$ with one $c$ for all $i$.

Now the clock. $z_{i}$ is $\tau'$-fixed and $R_{i}$ is fixed (ground) and invertible, so $\tau'_{t}(c) = c$. Take $k$ with $c \in B_{k}$; on $B_{k}$, $\tau'_{t} = \operatorname{Ad}(\exp(tH_{k} \otimes I_4))$ with $H_{k} = \sum_{j \le k}\omega'_{j}K_{j}$ real, so $c$ commutes with $H_{k}$. The $2K_{j}$ are commuting real involutions; in the irreducible real module of $B_{k}$, of dimension $2^{k}$, each of the $2^{k}$ joint sign patterns occurs exactly once (the $e_{j}$ flip the $j$th sign and fix the others). The least eigenvalue $-\tfrac12\sum_{j \le k}|\omega'_{j}|$ of $H_{k}$ is attained by one pattern only, because no $\omega'_{j}$ vanishes, so its spectral projection $P$, a real polynomial in $H_{k}$, has rank one. $\Gamma_{k} = \prod_{j \le k}e_{j}f_{j} \in B_{k}$ satisfies $\Gamma_{k}^{2} = 1$, commutes with $H_{k}$, and implements $\alpha$ on $B_{k}$; so $\Gamma_{k}P = \eta P$ with $\eta = \pm 1$ and $cP = \nu P$ with $\nu$ real, the range of $P$ being one real line preserved by both. Then $\alpha(c)P = \Gamma_{k}c\Gamma_{k}P = \eta^{2}\nu P = \nu P$ and $\mu P = c\,\alpha(c)P = \nu^{2}P$: $\mu = \nu^{2} > 0$.

So $z_{i}z_{j} = cR_{i}cR_{j} = c\,\alpha(c)R_{i}R_{j} = \mu R_{i}R_{j}$, and the elements $\hat R_{i} = \sqrt{\mu}\,R_{i} \in \Cl^{1}(p',q')$ satisfy $\hat R_{i}^{2} = z_{i}^{2} = x_{i}^{2}$ and anticommute pairwise. The assignment $x_{i} \mapsto \hat R_{i}$ is a grading-preserving homomorphism $\Cl(p,q) \to \Cl(p',q')$; it is injective because $\Psi(x_{B}) = \mu^{\lfloor |B|/2 \rfloor}c^{\,|B| \bmod 2}R_{B}$, so a relation among the $\hat R_{B}$ splits by parity into relations among the $\Psi(x_{B})$; and $n' = n$ by (ii) applied to the commutants, so it is an isomorphism of graded algebras. As in (ii), its even part gives $r' \equiv \pm r$, and for $r' \equiv -r$ with $r \not\equiv 0, 4$ the ungraded types differ, so $r' \equiv r$. Conversely the shifts of (ii) fix every plane and map ground into ground. Sharpness: $\Phi(x_{i}) = y_{i}e_{1}$, and $e_{1}$ commutes with every $K_{j}$, $j \ge 2$, and with $\omega'_{1}K_{1} = 0$; here $c = e_{1}$ has $c\,\alpha(c) = -1 < 0$, the sign the lowest-state argument excludes when no weight vanishes. If $\Psi$ intertwines clocks, $\Psi\tau_{t} = \tau'_{t}\Psi$, then $\Psi(x_{i}) = \Psi(\tau_{t}(x_{i})) = \tau'_{t}(\Psi(x_{i}))$, so the hypothesis holds.

\emph{(iv), half period.} 
On the floor $T_{k,\C}$ put $U_{t} = \exp\bigl(t\sum_{j \le k}\omega_{j}K_{j} \otimes I_4\bigr) = \prod_{j \le k}\bigl(\cos\tfrac{t\omega_{j}}{2} + \sin\tfrac{t\omega_{j}}{2}\,u_{j} \otimes I_4\bigr)$ (the scalar $1 \otimes I_4$ of the complexification, central by construction), $u_{j} = e_{j}f_{j}$, using $u_{j}^{2} = 1$ and that the $u_{j}$ commute. Each $u_{l}$ with $l \ne j$ commutes with $e_{j}, f_{j}$ and with every ground generator, and $u_{j}$ anticommutes with $e_{j}, f_{j}$; so $U_{t}$ commutes with the ground generators and with $K_{l}$ for $l > k$, which makes $\tau_{t} = \operatorname{Ad}(U_{t})$ consistent along the floors, and, writing $c = \cos\tfrac{t\omega_{j}}{2}$, $s = \sin\tfrac{t\omega_{j}}{2}$,
\[
\tau_{t}(e_{j}) = e_{j}(c - s\,u_{j} \otimes I_4)^{2} = \cos(t\omega_{j})\,e_{j} - \sin(t\omega_{j})\,f_{j} \otimes I_4, \qquad \tau_{t}(f_{j}) = \cos(t\omega_{j})\,f_{j} - \sin(t\omega_{j})\,e_{j} \otimes I_4,
\]
($e_{j}u_{j} = f_{j}$, $f_{j}u_{j} = e_{j}$). So $\tau$ fixes the ground algebra, $\tau_{t}$ is the identity on $P_{j}$ exactly when $t\omega_{j} \in 2\pi\mathbb{Z}$, and negates $e_{j}, f_{j}$ exactly when $t\omega_{j} \in \pi + 2\pi\mathbb{Z}$. The flow is the identity at $t$ if and only if $t\omega_{0}m_{j}/2\pi \in \mathbb{Z}$ for every $j$; since $\gcd_{j}m_{j} = 1$ some finite combination $\sum_{j}c_{j}m_{j} = 1$ with integers $c_{j}$, whence $t\omega_{0}/2\pi \in \mathbb{Z}$: the least period is $2\pi/\omega_{0}$. At $t = \pi/\omega_{0}$, $t\omega_{j} = \pi m_{j}$ with $m_{j}$ odd, so $\tau_{\pi/\omega_{0}}$ negates every plane generator: it is $\alpha_{B}$ on $B_{\infty}$, a real automorphism. For the common-rate boost $\kappa K_{\mathrm{tot}}$ of Theorem \claimref{synthesis.phase-flow-is-time-rotor-flow-kms-continuation}(iii), $K_{\mathrm{tot}} = \sum_{j}(N_{j} - \tfrac12)$, each $N_{j} - \tfrac12$ is the atom boost $\tfrac12\sigma_{j} = \tfrac12(v_{j}^{*}v_{j} - v_{j}v_{j}^{*})$, which for $v_{j} = \tfrac12(e_{j} + f_{j})$ and the frame $*$-structure ($e_{j}^{*} = e_{j}$, $f_{j}^{*} = -f_{j}$) is $\tfrac12 e_{j}f_{j} = K_{j}$; so all $m_{j} = 1$. (This reads $v_{j}$ as that theorem's annihilator, $N_{j} = v_{j}^{*}v_{j}$; in the tower's labelling, where $v_{j}$ creates, $N_{j} = v_{j}v_{j}^{*}$ and $N_{j} - \tfrac12 = -K_{j}$, so all $m_{j} = -1$, again odd, and nothing below changes.) If one plane has an even $m_{j}$, $\tau_{\pi/\omega_{0}}$ fixes $P_{j}$ and the half-period map is not the grading; the check shows both cases.

\emph{(v).} With $y_{i}^{2} = -x_{i}^{2}$, $(y_{i} \otimes I_4)^{2} = x_{i}^{2}$ and distinct images anticommute with each other and with the plane generators, so the $\C$-linear assignment extends to a homomorphism of the complexified limits with inverse $y_{i} \mapsto -x_{i} \otimes I_4$; it fixes every plane generator and hence $B_{\infty,\C}$, $\alpha_{B}$, each $P_{j,\C}$ and each $\tau$. It does not preserve the real forms.

The check verifies (i) by enumerating blades on floors $T_{1}, T_{2}$ against the planes of the next floor for ten ground signatures; the trace-form invariant of (ii) for $n \le 8$ and its invariance under the shift; that $\Phi$ sends $e_{2}, f_{1}$ to even elements while the shift fixes every plane; the framing lemma of (iii); the half-period and full-period identities of (iv) with odd and mixed-parity weights, and the lowest-state lemma of (iv): for weights with no zero the least eigenspace of $H_{k}$ has rank one ($2^{k}$ times the scalar part of its projection), $\Gamma_{k}$ acts on it by a sign, and $P\,c\,\alpha(c)\,P = \nu^{2}P$ for invariant $c$ in the commutant of $H_{k}$, while with one zero weight the rank is two and $e_{1}$ is invariant with $e_{1}\alpha(e_{1}) = -1$, in exact arithmetic with complex numbers as pairs of rationals.
